%!TEX program = lualatex
\documentclass[ngerman,headheight=70pt]{scrartcl}
\usepackage[ngerman]{babel}
\usepackage{blindtext}
\usepackage[utf8]{luainputenc}
\usepackage{scrlayer-scrpage}
\usepackage[T1]{fontenc}
\usepackage{fontspec}
\usepackage{multicol}
\usepackage{enumitem}
\usepackage[german=quotes]{csquotes}
\usepackage[bottom=2cm,footskip=8mm]{geometry}
\setmainfont{Source Sans Pro}
\MakeOuterQuote{"}

\setlength{\parindent}{0pt}
\setlength{\parskip}{6pt}

\pagestyle{plain.scrheadings}
\thispagestyle{scrheadings}
\lohead{\vspace*{-1.5cm}\includegraphics[scale=0.4]{../images/stupa_logo_rot.pdf}}
\rohead{\\Präsidium des\\Studierendenparlaments}

\lofoot[Universität Hamburg · Präsidium des Studierendenparlaments\\\hrule Von-Melle-Park 5 · D-20146 Hamburg]{Universität Hamburg · Präsidium des Studierendenparlaments\\\hrule Von-Melle-Park 5 · D-20146 Hamburg}
\cofoot[]{}
\cohead[\thepage]{}

\renewcommand{\thesubsection}{TOP \arabic{subsection}}
\renewcommand{\thesubsubsection}{\arabic{subsubsection}.}
\setcounter{subsection}{-1}

\begin{document}
    UHH · StuPa-Präsidium · Von-Melle-Park 5 · D-20146 Hamburg

    %%% PAGE ONE (EDIT LATER) %%%%
    % Hier den Titel des Protokolls eintragen
    \section*{Protokoll der 9. Sitzung des Studierendenparlaments vom 10. November 2016}

    % Protokollanten eintragen
    \textbf{Protokoll: Jim Martens}\\
    \textbf{Ort: VMP5 0079}\\
    \textbf{Beginn: 18.29 Uhr}\\
    \textbf{Ende: 23.55 Uhr}

    % Wohl die komplizierteste Stelle im Protokoll
    % Als Faustregel sollten nicht mehr als 3 Namen auf einer Zeile stehen
    % Einfach an dieser Vorlage orientieren
    % Im Idealfall ist hier nicht viel zu ändern, da kaum jemand fehlt.
    % Rücktritte müssen natürlich pro Protokoll angepasst werden
    \vspace{0.5cm}
    \begin{tabular}{ll}
        Anwesend: & \\
            RCDS (5 Sitze): &  \\
             CampusGrün (14 Sitze): & Laura Franzen, Geoffrey Youett, Elena Rysikova, \\
                                   & Philipp Droll, Yasemin Günther, Melf Johannsen, \\
                                   & Mario Moldenhauer, Christian Lagod, Jim Martens, \\
                                   & Svenja Horn, Martin Sievert, Armin Günther \\
             BGZM-Liste (2 Sitze) : & \\
             WiWi (2 Sitze): & \\
             Unicorns (5 Sitze): &  Johannes Peplow, Annkathrin Löffler \\
             Liste LINKS (3 Sitze): & Gunhild Berdal, Till Petersen, Sinah Mielich \\
             HWP (2 Sitze): &  \\
             MIN (4 Sitze): & Ailina Salten \\
             SDS* (3 Sitze): & Mena Winkler, Jacob Petersein, Artur Brückmann \\
             Bart-LISTE (2 Sitze): & Timo Zeimet, Dominic Laumer \\
             LHG (1 Sitz): & Tobias Heisig \\
             harte zeiten (1 Sitz): & Tobias Berking \\
             Jura (1 Sitz): & \\
             AL (2 Sitze): & Karima Schulze, Henri Weber \\
            & \\
        Entschuldigt: & Ramon Weilinger, Antonia Niecke, Ramin Shakiba, \\
                      & Jennifer Maack, Benjamin Welling (RCDS), Jakob Pape, \\
                      & Karen Martirosian (BGZM), Charleen Lorenz (WiWi), Katharina Kucza, \\
                      & Andreas Hartkamp, Marielle Hermstrüwer (Unicorns), \\
                      & Jan Detampel (MIN) \\
                      &\\
        Unentschuldigt abwesend: & Tahnee Herzig (CG), Claas-Friso Hente (WiWi), Ajdina Karahasan,\\
                                & Gökce Celik (HWP), Lotte Rullkötter, Nabil Jabareen (MIN), \\
                                & Maurice Martens (Jura) \\
                                &\\
        Rücktritte: & \\
        Mandatsverluste: & Kay Zöllmer \rightarrow Charleen Lorenz (WiWi), \\
                         & Samet Gunay \rightarrow Gökce Celik (HWP), \\
                         & Johann Baumhöfener \rightarrow Maurice Martens (Jura) \\
    \end{tabular}
    \newpage
    %%% Vorgeschlagene Tagesordnung %%%
    \underline{Vorgeschlagene Tagesordnung}
    \begin{enumerate}[label={\textbf{Top \theenumi}},leftmargin=*]
        % folgende Zeile bei Sitzungsfortsetzungen einkommentieren und x durch
        % die Zahl des zuletzt behandelten TOPs ersetzen
        %\setcounter{enumi}{x}
        \item Debatte zu G20 (V1617-041) (60 Min.)
        \item Vorbereitung der Wahlen zum Studierendenparlament (60 Min.)
        \begin{enumerate}
            \item Diskussion über Konzeptvorschlag des Präsidiums (V1617-034)
            \item Diskussion über die Steigerung der Wahlbeteiligung
            \item Beschlussfassung über den Zeitplan und den Urnenplan
            (V1617-035, V1617-036)
        \end{enumerate}
        \item Debatte über den Haushalt der Freien und Hansestadt Hamburg
              2017/2018 (V1617-038) (60 Min.)
        \item Aufruf zur Mahnwache zum 78. Jahrestag der Pogromnacht (V1617-037) (30 Min.)
        \item Erhöhung Wochenarbeitszeit für Festanstellung im Finanzreferat des AStA (V1617-039) (15 Min.)
        \item Kaffeeautomat im AStA-Infotreff (V1617-040) (15 Min.)
        \item Verschiedenes (5 Min.)
    \end{enumerate}

    \newpage

    % Hier beginnt Teil A. Bei Fortsetzungen kommt hier nur der Teil über
    % die Feststellung der endgültigen Fassung der Tagesordnung.
    % Für sonstige Sitzungen ist dies bereits das Gerüst für die Formalia.
    % Man muss nur noch die protokollierten Inhalte eintragen.

    {\Large \underline{Teil A}}

    \subsection{Formalia}

    \textit{Gunhild eröffnet die Sitzung um 18.35 Uhr.}

    Tobias Heisig erklärt sich bereit die Redeliste zu führen in Abwesenheit
    von Ramon Weilinger.

    \subsubsection{Geschäftsbericht Präsidium}

    Gunhild berichtet aus dem Präsidium.

    \begin{itemize}
        \item Vertrag mit der Druckerei wurde abgeschlossen
        \item Mandatsverlustregelung wird von Gunhild erklärt und Näheres zum Entschuldigen
      gesagt.
        \item Gunhild erwähnt die inhaltliche Veranstaltung, die in 2 Wochen stattfindet.
        \item Jim erwähnt Veränderungen an den Kandidaturbögen und erinnert an die Fristen
      zur Abgabe von den Bögen.
        \item Gunhild erwähnt, dass wir über die US-Wahlen geredet haben.
    \end{itemize}

    Gunhild nennt die Personen, die ihr Mandat verloren haben und wer dafür nachrückt.

    \subsubsection{Anfragen an das Präsidium}

    F: Till fragt, ob die nachgerückten Leute eingeladen wurden. \\
    A: Ja.

    F: Karima fragt, was die Relevanz der US-Wahl spezifisch für das StuPa
    ist. \\
    A: Gunhild antwortet mit ihrer persönlichen Meinung.

    \subsubsection{Geschäftsbericht AStA}

    Franzi berichtet aus dem AStA. Der AStA hat sich weiter mit dem G20-Gipfel
    beschäftigt. Es gab ein Treffen zur Schulpolitik in Hamburg. Vielleicht wird
    ein Kongress für inklusive Bildung in Hamburg stattfinden. Es gab Veranstaltungen
    zum Gedenken an die Pogromnacht mit Esther Bejarano am 3.11. und eine Mahnwache
    am 9.11.

    Philipp fährt mit dem Bericht vor. Es gab letzte Woche ein Treffen mit Herrn
    Hecht und dem Architekten der Universität. Es wird sich mit der baulichen
    Entwicklung des Unicampus in den nächsten 50 Jahren beschäftigt. Die Uni liegt
    mitten in der Stadt, ist aber kaum mit ihr verbunden. Es gibt am 15.11. ein
    von 13-20 Uhr dauerndes Symposium zu dem Thema.

    Timo berichtet, dass der Jahresabschluss noch nicht ganz durch ist und in den
    nächsten Wochen kommen soll.

    Geoffrey war auf einem Freiraumtreffen und berichtet davon. Es gab außerdem
    am Mittwoch anlässlich des US-Wahlergebnisses ein Treffen mit rund 20 Leuten
    über die Konsequenzen der Wahl.

    Karima berichtet über die Organisation der Politkneipe. Ebenfalls erzählt sie
    über die neuen Gleichstellungsrichtlinien der Uni. Es ist geplant ein Filmabend
    mit dem Film "Triumph des guten Gewissens" zu veranstalten.
    Mario ergänzt, dass die Geschlechterangabe bereits im nächsten Semester auf
    den Semesterunterlagen (Studierendenausweis, Semesterticket) entfallen
    soll.

    \subsubsection{Anfragen an den AStA}

    Tobias Berking fragt Philipp über das Treffen zum Unibau (Sinn des Treffens
    und Hintergrund der Einschätzungen) und ergänzt die Frage
    mit inhaltlichen Anmerkungen.

    \textit{Zeitansage: 19.10 Uhr}

    Philipp antwortet, dass es keinen konkreten Anhaltspunkt für das Treffen gab.
    Stattdessen soll es um die gesamtuniversitäre bauliche Entwicklung gehen.
    Dabei geht es weniger um Schönheit der Gebäude, sondern um Funktionalität.

    Tobias weist auf die geschichtliche Relevanz der Gebäude hin.

    Till weist darauf hin, dass AStA-Leute vorbereitet zum Symposium gehen sollten.
    Rückbau von Gebäuden kann viel erreichen. Transparenz ist ein wichtiger Punkt.

    Franzi erwähnt, dass die Geschichte der Universität und des Grindelviertels
    diskutiert wurde auf dem Treffen. Mit den Überlegungen wird sich von der
    Beschränkung der Schuldenbremse gelöst und versucht Funktionalität, Ästhetik
    und Geschichte miteinander zu verbinden.

    \subsubsection{Dringlichkeitsanträge des AStA}

    Keine.

    \subsubsection{Aktuelle Stunde zum US-Wahlergebnis}

    Jim stellt einen Antrag eine Aktuelle Stunde zum Ausgang der US-Wahl zu machen.

    Es gibt mehrere Redebeiträge.

    \textit{Zeitansage: 20.28 Uhr}

    \subsubsection{Feststellung der endgültigen Fassung des Teils B der Tagesordnung}

    Vorschlag vom Präsidium ein TOP zur Bestätigung der stellv. Referent*in des
    Alle*Frauen-Referats einzufügen.
    Es gibt keinen Widerspruch gegen die Aufnahme in die Tagesordnung.
    Es wird mehrheitlich für Position 1 gestimmt.

    Die AL stellt einen Änderungsantrag zur Tagesordnung. Aktuelle TOPs 5 und 6
    sollen vor der Debatte zu G20 kommen. Dies wird mehrheitlich abgelehnt.

    Tagesordnung wird mehrheitlich bei 3 Gegenstimmen und 2 Enthaltungen angenommen.

    %%% BESCHLOSSENE TAGESORDNUNG %%%
    \underline{Endgültige Fassung der Tagesordnung}
    \begin{enumerate}[label={\textbf{Top \theenumi}},leftmargin=*]
        % folgende Zeile bei Sitzungsfortsetzungen einkommentieren und x durch
        % die Zahl des zuletzt behandelten TOPs ersetzen
        %\setcounter{enumi}{x}
        \item Bestätigung der stellvertretenden Referent*in des Alle*Frauen-Referates
        \item Debatte zu G20 (V1617-041) (60 Min.)
        \item Vorbereitung der Wahlen zum Studierendenparlament (60 Min.)
        \begin{enumerate}
            \item Diskussion über Konzeptvorschlag des Präsidiums (V1617-034)
            \item Diskussion über die Steigerung der Wahlbeteiligung
            \item Beschlussfassung über den Zeitplan und den Urnenplan
            (V1617-035, V1617-036)
        \end{enumerate}
        \item Debatte über den Haushalt der Freien und Hansestadt Hamburg
              2017/2018 (V1617-038) (60 Min.)
        \item Aufruf zur Mahnwache zum 78. Jahrestag der Pogromnacht (V1617-037) (30 Min.)
        \item Erhöhung Wochenarbeitszeit für Festanstellung im Finanzreferat des AStA (V1617-039) (15 Min.)
        \item Kaffeeautomat im AStA-Infotreff (V1617-040) (15 Min.)
        \item Verschiedenes (5 Min.)
    \end{enumerate}

    \subsubsection{Feststellung der Beschlussfähigkeit}

    \textit{Mit 27 Parlamentarier*innen ist das Parlament beschlussfähig.}

    \subsubsection{Genehmigung der Protokolle der vorangegangenen Sitzungen}

    Till merkt an, dass bei CampusGrün in der letzten Sitzung insgesamt nur 13
    Personen in der Anwesenheitsliste stehen. Das Präsidium wird dies prüfen
    und falls nötig korrigieren.

    Das Protokoll wird mehrheitlich bei 0 Gegenstimmen und 3 Enthaltungen
    angenommen.

    \textit{Präsidiumspause bis 20.55 Uhr (dynamisch verlängert bis 21.05)}

    \vspace{0.5cm}
    {\Large \underline{Teil B}}

    \subsection{Bestätigung der stellvertretenden Referent*in des Alle*Frauen-Referates}

    Lisa Stamm wurde als neue stellvertretende Referent*in des Alle*Frauen-Referates
    vom Sprecher*innenrat gewählt.

    Lisa Stamm wird einstimmig bestätigt.

    \subsection{Debatte zu G20}

    Franziska stellt den Antrag unter Berücksichtigung der aktuellen Entwicklungen
    vor. Am Sonntag trifft sich die AG, die den Alternativgipfel zu G20 organisiert.

    Armin sagt, der Fokus liege auf der globalen Bedeutung von G20 statt den konkreten
    Auswirkungen auf Hamburg.

    Franziska berichtet vom Treffen des Jugendbündnisses in dieser Woche.

    Elias ergänzt u.a. Infos zum Potsdamer Abkommen und dessen geschichtlicher
    Relevanz.

    \textit{Zeitansage: 21.17 Uhr}

    Gunhild mahnt an, dass man sich angesichts von Trump frühzeitig engagieren
    muss, um den Gipfel zu verhindern. Es sollte herausgestellt werden, dass
    G20 eine Gegenorganisation zur UN und keine Friedensstiftende ist.

    Johannes stellt den Änderungsantrag in Zeile 25 "japanischen" zu streichen.
    Es gibt keinen Widerspruch.

    Tobias Berking weist darauf hin, dass die G20 in der Tat eine Gegenorganisation
    sind.

    Till stellt den Antrag in Zeile 103 nach dem Wort "Protestbewegung" folgendes
    einzufügen.

    \begin{blockquote}
        Dafür ist es erforderlich bei allen Aktivitäten - einschließlich den hier
        zum Schluss des Antrags genannten - zu erarbeiten und deutlich zu machen,
        was die fortschrittliche, friedenspolitische und soziale Alternative zu
        den von den G20 repräsentierten Positionen ist. Diese Alternative muss
        durch die Bevölkerungen international erwirkt werden.
    \end{blockquote}

    Der Antrag wird mehrheitlich mit einer Gegenstimme und einer Enthaltung angenommen.

    \subsection{Vorbereitung der Wahlen zum Studierendenparlament}

    \subsubsection{Diskussion über Konzeptvorschlag des Präsidiums}

    Jim stellt das inhaltliche Konzept vor.

    Gunhild ergänzt, dass sich das Präsidium bezüglich des Konzepts einig ist.

    Philipp begrüßt das Konzept und appelliert an alle einen politischen Wahlkampf
    zu machen.

    Gunhild findet das veränderte Konzept der Wahlveranstaltungen besser als was
    bisherige Präsidien gemacht haben.

    Till erinnert an einen Antrag zu einem fairen Wahlkampf und fragt, ob an einen
    solchen Antrag bereits gedacht wurde. Till schlägt vor selber eine Publikation
    seitens des Präsidiums zu machen, um das Wahlmotto mit Inhalt zu füllen. Er
    schlägt ebenfalls vor allen Listen die Möglichkeit zu geben auf vielleicht 2
    Zeilen vorzustellen, wie sie die Universität entwickeln möchten.

    Jim hat die Idee eine Publikation vor der Briefwahl zu machen und dort auf
    die Wahl, die Termine und die Listen hinzuweisen.

    Elias schlägt vor im Wahlmotto "selbstgestaltet" durch "selbstgestaltend"
    zu ersetzen.

    Gunhild redet darüber, dass aktuell viele Gespräche stattfinden und sich überlegt
    werden muss, wie man Leute dazu bekommt sich aktiv einzubringen.

    Timo Zeimet fragt, was Gunhild mit aktiv meint. Gunhild antwortet, dass aktiv
    sich auf alles bezieht.

    Till findet eine klare Positionierung und eine Authentizität für wichtig.
    Allerdings ist es nicht so toll, wenn man sich gar nicht auf andere bezieht.
    Es muss zwischen einer Diffamierung und einer authentischen Positierung
    unterschieden werden.

    Timo Zeimet stellt klar, dass er Schlammschlachten im Wahlkampf auch schlecht
    findet, eine inhaltliche Positionierung jedoch nicht.

    Jim äußert seine Hoffnung, dass die Listen sich in diesem Wahlkampf inhaltlicher
    präsentieren.

    \subsubsection{Diskussion über die Steigerung der Wahlbeteiligung}

    Gunhild liest die Frage von Victoria Mader vor (siehe Anhang).

    Geoffrey weist darauf hin, dass viele der Parlamentarier*innen, die letzte Sitzung
    fehlten, auch diese Sitzung fehlen und die Frage daher an diese weitergeleitet
    werden sollte. Das Präsidium wird dies tun.

    Timo Zeimet merkt an, dass man vielleicht dem "Otto-Normalstudierenden"
    besser erklären können sollte, was das StuPa konkret mit deren Studium zu tun
    hat.

    Franziska entgegnet, dass z.B. eine Studie gemacht wurde in Baden-Württemberg,
    warum Studierende in die Gewerkschaft verdi eingetreten sind. Ergebnis war, dass
    es hauptsächlich aufgrund der Solidarität und weniger wegen der konkreten
    Auswirkungen auf Praktika geschah.

    Timo Zeimet sagt, dass vor einer Solidarität eine Identifikation stehen muss.
    Diese Identifikation fehle jedoch größtenteils.

    Laura ist der Überzeugung, dass die gedankliche Verknüpfung zwischen StuPa
    und dem eigenen Studium defacto kaum vorhanden ist. Diese Verknüpfung müsste
    besser herausgestellt werden, damit Studierende sich damit identifizieren können.

    Tobias Berking merkt an, dass einem das Ba/Ma-System dauerhaft nahelegt
    egoistisch zu sein. Dagegen muss man argumentieren. Themen wie G20 gehen sehr
    wohl alle Studierenden an, da es um die Entwicklung der Menschheit geht, deren
    Teil sie sind.

    Till weist darauf hin, dass es in ganz Deutschland unterschiedliche StuPas gibt,
    deren Wahlbeteiligungen grundsätzlich bei maximal 20\% stehen. Früher war dies
    jedoch anders und es haben über die Hälfte der Studierenden gewählt.
    Der Kernpunkt ist, dass sich die Menschen früher ernster genommen haben als
    politischen Faktor. Heute rechnen sie sich kaum eine Bedeutung aus.

    Ein Gast der BART-Liste merkt an, dass unpolitischer Wahlkampf erfolgreich war
    und den Menschen geholfen hat, überhaupt erst einmal die politische Meinung
    zu äußern. Diese seien häufig bezüglich ihrer politischen Position verunsichert.

    Franziska redet darüber, was die wirkliche Ursache der Verunsicherung ist.

    Timo Zeimet ist der Meinung, dass den Studierenden bewusst gemacht werden muss,
    dass u.a. das StuPa da und relevant ist.

    Laura weist darauf hin, dass ein Zusammenhang zwischer Fachkultur und
    gesamtuniversitärer Ebene geschaffen werden muss. Auch seien viele davon überzeugt
    ohnehin wenig erreichen zu können. Neuen Menschen sind häufig die personellen
    Beziehungen nicht bekannt und dies kann sehr einschüchternd wirken.
    Diesbezüglich sei eine bessere Informationspolitik nötig.

    Gunhild sagt, dass fast alle Präsidien sich bemühen den Zusammenhang zwischen
    StuPa und dem Studienalltag klarzumachen. Sie ist der Überzeugung, dass
    mehr Politik besser überzeugt als unpolitisches Handeln.

    \textit{Zeitansage: 22.40 Uhr}

    Till ist der Überzeugung, dass Leute Gesprächsbedarf haben und daher ins
    Gespräch kommen, wenn sie angesprochen werden.

    Ein Gast hält einen Redebeitrag.

    Timo Zeimet sagt, dass Politik verständlich sein sollte für den Leser.

    Mario fasst zusammen, dass er zwei Lager sieht. Er stellt die Frage, ob
    die Strategie des "linken" Blocks die richtige ist, wenn sie seit Jahren die
    gleiche ist und sich trotzdem die Wahlbeteiligung nicht erhöht.

    Till ist der Meinung, dass es darum geht, was richtig und falsch ist. Die Position
    der Linkspartei zu Bundeswehreinsätzen beispielsweise sei korrekt, auch wenn
    sie nicht die Mehrheit der Stimmen hat.

    \subsubsection{Beschlussfassung über den Zeitplan und den Urnenplan}

    Wir kommen zum Zeitplan. Geoffrey fragt, ob das Präsidium vorhätte die Fristen
    laxer zu handhaben. Gunhild antwortet, dass die Fristen fest eingehalten werden.

    Der Zeitplan wird mehrheitlich bei einer Gegenstimme und einer Enthaltung
    angenommen.

    Wir kommen zum Urnenplan. Till beantragt Urne 6 zu streichen. Es gibt keinen
    Widerspruch.

    \textit{Geoffrey beantragt Verlängerung der Sitzung bis nach TOP 8. Der Antrag
    wird mehrheitlich bei 2 Gegenstimmen und 3 Enthaltungen angenommen.}

    Franziska weist darauf hin, dass in der Chemie so gut wie nichts im Foyer
    möglich ist. Das Gleiche gilt auch für die HWP. Dies kann ebenso den
    Pferdestall und das Rechtshaus betreffen.

    Till überlegt, ob es Sinn macht die Urne der Holzwirtschaft wieder zu öffnen.

    Despina weist darauf hin, dass in den Naturwissenschaften viele Leute nicht
    zum Campus kommen.

    Timo ist besorgt, dass eine neue Urne in der HoWi nicht besetzt werden kann.

    Till schlägt vor eine 19. Urne an der Holzwirtschaft (nach Rücksprache mit
    dem FSR) einzurichten. Es gibt keinen Widerspruch.

    Der Urnenplan wird angenommen.

    \textit{Präsidiumspause bis 23:20, dynamisch bis 23:23}

    \subsection{Debatte über den Haushalt der Freien und Hansestadt Hamburg}

    Jim stellt seinen Antrag vor.

    Philipp stellt den Antrag des AStA vor.
    Am 6.12. soll eine Vollversammlung zum Thema stattfinden.

    Franziska ergänzt ihre Ansichten. Sie stellt mehrere Änderungsanträge an dem
    Antrag des AStA vor, welche von den Antragsstellern übernommen werden.

    Der Punkt hinter "Kritik" in Zeile 15 wird in einen Doppelpunkt geändert.
    Direkt danach wird "Mit der Schuldenbremse soll der Eindruck erweckt werden,"
    eingefügt. "Es fehlt" wird geändert in "es fehle" und "nur viel zu oft"
    wird gestrichen.

    In Zeile 21 wird "der gesellschaftlichen Missstände" ersetzt durch "für eine
    humane Gesellschaftsentwicklung". In Zeile 22 wird "autonome" durch
    "gesellschaftlich verantwortliche" ersetzt.

    Nach Zeile 41 wird folgendes eingefügt.

    \begin{blockquote}
        Das Studierendenparlament beschließt zur Diskussion und Positionierung
        zu dieser Auseinandersetzung am 6. Dezember eine Vollversammlung
        einzuberufen.
    \end{blockquote}

    Bei der Abstimmung wird zuerst Jims Antrag und anschließend der vom AStA
    abgestimmt.

    Es ergeben sich 7 Stimmen für und 0 Stimmen gegen Jims Antrag bei einer
    Mehrzahl an Enthaltungen. Damit ist der Antrag beschlossen.

    Für den Antrag des AStA gibt es 11 Stimmen, bei einer Gegenstimme und vielen
    Enthaltungen. Damit ist auch der Antrag beschlossen.

    \subsection{Aufruf zur Mahnwache zum 78. Jahrestag der Pogromnacht}

    Till stellt einen Änderungsantrag vor. Es gibt keinen Widerspruch.

    Es soll ab Zeile 1 "Das SP möge" bis Zeile 2 "anschließen" durch folgenden
    Text ersetzt werden.

    \begin{blockquote}
        Das StuPa begrüßt den Beschluss des Akademischen Senats vom 13. Oktober
        2016 zum Aufruf zur Mahnwache zum 78. Jahrestag der Pogromnacht und das
        Bemühen, daraus einen gemeinsamen Aufruf von StuPa und AS zu machen.
        Das StuPa ist bestrebt, dass zur Mahnwache im nächsten Jahr ein gemeinsamer
        Aufruf gelingt.
    \end{blockquote}

    Der Antrag wird entsprechend angenommen.

    \subsection{Erhöhung Wochenarbeitszeit für Festanstellung im Finanzreferat des AStA}

    Franziska stellt den Antrag vor.

    Till stellt die Frage, ob in den erhöhten Arbeitsstunden vorgesehen ist,
    dass die Festangestellten an den AStA-Sitzungen teilnehmen.

    Franziska antwortet, dass es ein wöchentliches Treffen gibt zwischen AStA-Vorstand
    und Festangestellten, aber die höheren Arbeitsstunden inkludieren nicht
    die Teilnahme an den AStA-Sitzungen.

    Der Antrag wird einstimmig angenommen.

    \subsection{Kaffeeautomat im AStA-Infotreff}

    Philipp stellt den Antrag vor.

    Franziska erklärt, warum es den Antrag gibt.

    Christian weist darauf hin, dass der Wirtschaftsrat dies bereits vorläufig
    beschlossen hat.

    Es gibt 13 Stimmen für den Antrag bei 3 Gegenstimmen und 7 Enthaltungen.
    Damit ist der Antrag angenommen.

    \subsection{Verschiedenes}

    \textit{Gunhild schließt die Sitzung um 23.55 Uhr.}

\end{document}
